\documentclass[12pt]{article}
\usepackage[utf8]{inputenc}
\usepackage[T1]{fontenc}
\usepackage{hyperref}
\title{Saterfriesisch Language Resources for K-12 Curriculum Development}
\author{Compiled by Claude Sonnet 5.5 (claude-sonnet-5-5) for Andr\'as Kornai}
\date{}
\begin{document}
\maketitle

\section{Introduction}
Saterfriesisch (Saterland Frisian, ISO 639-3 stq; Glottocode sate1242) is an East Frisian language spoken in the Saterland in Lower Saxony, Germany. This report rests on a small number of web searches. Sources give about 2,000 native speakers (2015 figure), while the survey reference gives about 1,000. UNESCO classes the language as definitely endangered.

\section{Primary Educational Resources}
The searches found no confirmed curriculum, textbook or teaching hours. They also found no information on teacher training or on whether Saterfriesisch is a regular school subject.
\begin{itemize}
\item Schulzentrum Saterland in Lower Ramsloh, a general school with primary and secondary levels, founded in 1971. A description of the school appears on Wikipedia. The search did not confirm what Saterfriesisch teaching it offers. \url{https://en.wikipedia.org/wiki/Schulzentrum_Saterland}
\item Mercator Research wiki entry on Saterfrisian in Germany. It says new vocabulary is introduced by Saterland schools, associations and work groups. \url{https://wiki.mercator-research.eu/languages:saterfrisian_in_germany}
\end{itemize}

\section{Linguistic and Reference Materials}
\begin{itemize}
\item Marron C.\ Fort, \emph{Saterfriesisches W\"orterbuch}, first published in 1980. A second, revised and greatly expanded edition appeared in 2015 (Helmut Buske Verlag, ISBN 978-3-87548-723-7). The publisher record found: \url{https://www.lehmanns.ch/media/87178986}
\item Fryske Akademy and Oldenburgische Landschaft: a Saterland Frisian grammar (Sprachlehre, 2022) and an English-language scholarly online grammar in progress. \url{https://fryske-akademy.nl/en/news-and-events/news-item/news/detail/lectures-for-the-oldenburgische-landschaft-in-saterland}
\item Overview of the language (Wikipedia). \url{https://en.wikipedia.org/wiki/Saterland_Frisian_language}
\item Academic article listed at De Gruyter Brill. \url{https://www.degruyterbrill.com/document/doi/10.1075/nss.33.pi/html}
\end{itemize}

\section{Religious and Community Resources}
No religious texts were identified. Language work in the region is linked to the Oldenburgische Landschaft and the Seeltersk Kontoor, mentioned in the sources above.

\section{Digital and Audio Resources}
\begin{itemize}
\item Fryske Akademy project on speech technology (``speechology'') for Saterland Frisian. \url{https://fryske-akademy.nl/en/news-and-events/news-item/news/detail/speechology-for-the-frisian-language-of-saterland}
\end{itemize}

\section{Curriculum Development Context}
The language has a published dictionary, a grammar and active academic interest. Verified information on classroom teaching was not found. Curriculum planning should start from the Saterland school system and the language bodies named above.

\end{document}